\section{研究思路}

\begin{frame}{研究思路与主要贡献}
  \begin{columns}[T,onlytextwidth]
    \column{0.52\textwidth}
    \begin{block}{研究思路}
      \begin{itemize}
        \item 从 \hlb{动态评估框架} 出发重构信用风险预警任务
        \item 用 LSTM 建模企业连续财务轨迹
        \item 用 Logistic 回归提供静态基准与可解释参照
        \item 在实证比较基础上设计混合风控方案
      \end{itemize}
    \end{block}

    \column{0.44\textwidth}
    \begin{exampleblock}{主要贡献}
      \begin{enumerate}
        \item 构建可操作的动态评估框架
        \item 证明时序信息对风险辨别有增益
        \item 识别两类模型的优势边界
        \item 提出两阶段混合风控体系
      \end{enumerate}
    \end{exampleblock}
  \end{columns}
\end{frame}

\begin{frame}{技术路线：从面板数据到业务预警}
  \centering
  \begin{tikzpicture}[
    node distance=0.55cm,
    box/.style={draw=CsuBlue,rounded corners,fill=CsuBlue!6,minimum width=2.25cm,minimum height=0.95cm,align=center,font=\small},
    arr/.style={-{Stealth[length=2.2mm]},thick,CsuBlue}
  ]
    \node[box] (data) {A股上市公司\\2007--2022};
    \node[box,right=of data] (feature) {特征构建\\标准化处理};
    \node[box,right=of feature] (slice) {时序切片\\训练/验证/测试};
    \node[box,right=of slice] (model) {Logistic\\LSTM};
    \node[box,right=of model] (eval) {AUC / KS\\Precision / Recall};
    \node[box,below=0.9cm of model] (hybrid) {混合模型\\初筛 + 复核};

    \draw[arr] (data) -- (feature);
    \draw[arr] (feature) -- (slice);
    \draw[arr] (slice) -- (model);
    \draw[arr] (model) -- (eval);
    \draw[arr] (model) -- (hybrid);
    \draw[arr] (hybrid) -| (eval);
  \end{tikzpicture}

  \vspace{0.8em}
  \begin{itemize}
    \item 研究不只比较“谁更准”，更关注 \hl{哪类模型适合哪类风控任务}。
    \item 最终目标是把模型比较结果转化为 \hlb{可落地的业务流程设计}。
  \end{itemize}
\end{frame}
